\documentclass{beamer}
\usepackage[utf8]{inputenc}
\usepackage[spanish]{babel}
\usepackage{listings}
\usepackage{tikz}
\usepackage{color}

% Colors
\definecolor{myblue}{RGB}{0, 0, 150} % Dark blue for text
\setbeamercolor{background canvas}{bg=white}
\setbeamercolor{normal text}{fg=myblue}
\setbeamercolor{frametitle}{fg=myblue}
\setbeamercolor{title}{fg=myblue}
\setbeamercolor{itemize item}{fg=myblue}
\setbeamercolor{itemize subitem}{fg=myblue}

% Listings config for Java
\lstset{
  language=Java,
  basicstyle=\ttfamily\footnotesize\color{myblue},
  keywordstyle=\bfseries\color{blue},
  stringstyle=\color{cyan},
  commentstyle=\color{gray},
  showstringspaces=false
}

\title{Ayudantía [7 Oct]}
\date{}
\author{}

\begin{document}

\begin{frame}
  \titlepage
\end{frame}

% --- Section: Ordenamiento en Java ---
\section{Librería Estándar de Java para Ordenar}

\begin{frame}{¿Qué es el ordenamiento en Java?}
  \begin{itemize}
      \item Java provee interfaces estándar para definir cómo se deben comparar y ordenar los objetos.
      \item Las principales son \texttt{Comparable}, \texttt{Comparator} y \texttt{Collator}.
      \item Permiten usar métodos como \texttt{Collections.sort()} o \texttt{Arrays.sort()} fácilmente.
  \end{itemize}
  \vspace{0.5cm}
  \centering
  \begin{tikzpicture}
    \node[draw, rounded corners, align=center, text=myblue, draw=myblue] at (0,0) {Objetos \\ Desordenados};
    \draw[->, thick, myblue] (1.5,0) -- (3,0) node[midway, above] {\footnotesize \texttt{sort()}};
    \node[draw, rounded corners, align=center, text=myblue, draw=myblue] at (4.5,0) {Objetos \\ Ordenados};
  \end{tikzpicture}
\end{frame}

\begin{frame}{\texttt{Comparable} y \texttt{Comparator}: Concepto}
  \textbf{Comparable:}
  \begin{itemize}
      \item Define el orden "natural" de los objetos de una clase.
      \item La clase misma debe implementar la interfaz.
  \end{itemize}
  \vspace{0.3cm}
  \textbf{Comparator:}
  \begin{itemize}
      \item Define un orden externo.
      \item Útil para ordenar de múltiples formas o clases que no podemos modificar.
  \end{itemize}
  \vspace{0.3cm}
  \textbf{Collator:}
  \begin{itemize}
      \item Especializado en Strings. Considera reglas específicas de idioma (acentos, ñ, etc).
  \end{itemize}
\end{frame}

\begin{frame}[fragile]{Definición: \texttt{Comparable} (Java)}
\begin{lstlisting}
public interface Comparable<T> {
    /**
     * Compara este objeto con el objeto especificado.
     * Retorna:
     * Un entero negativo, cero, o un entero positivo 
     * si este objeto es menor, igual, o mayor que 
     * el objeto especificado.
     */
    public int compareTo(T o);
}
\end{lstlisting}
\end{frame}

\begin{frame}[fragile]{Definición: \texttt{Comparator} (Java)}
\begin{lstlisting}
public interface Comparator<T> {
    /**
     * Compara sus dos argumentos para el orden.
     * Retorna:
     * Un entero negativo, cero, o un entero positivo
     * si el primer argumento es menor, igual, o mayor 
     * que el segundo.
     */
    int compare(T o1, T o2);
}
\end{lstlisting}
\end{frame}


% --- Section: Árboles Binarios ---
\section{Árboles Binarios}

\begin{frame}[fragile]{Árboles Binarios: Definición}
  Un árbol binario es una estructura de datos jerárquica no lineal.
  \begin{itemize}
      \item Está compuesto de \textbf{nodos}.
      \item Cada nodo contiene un valor y referencias a como máximo dos hijos: el \textbf{hijo izquierdo} y el \textbf{hijo derecho}.
      \item El nodo superior se llama \textbf{raíz}.
      \item Los nodos sin hijos se llaman \textbf{hojas}.
  \end{itemize}
  \vspace{0.3cm}
  \centering
  \begin{tikzpicture}[level/.style={sibling distance=30mm/#1}, every node/.style={text=myblue, draw=myblue}]
    \node [circle,draw] (z){A}
      child {node [circle,draw] (a) {B}
        child {node [circle,draw] (b) {D}}
        child {node [circle,draw] (g) {E}}
      }
      child {node [circle,draw] (j) {C}
        child {node [circle,draw] (k) {F}}
        child[missing] {}
      };
  \end{tikzpicture}
\end{frame}

\begin{frame}{Árbol Binario de Búsqueda (BST)}
  Un tipo especial de árbol binario que facilita la búsqueda rápida.
  \vspace{0.3cm}\\
  \textbf{Propiedades:}
  \begin{itemize}
      \item El subárbol izquierdo de un nodo solo contiene nodos con claves \textbf{menores} que la clave del nodo.
      \item El subárbol derecho de un nodo solo contiene nodos con claves \textbf{mayores} que la clave del nodo.
      \item Tanto el subárbol izquierdo como el derecho deben ser también árboles binarios de búsqueda.
  \end{itemize}
\end{frame}

\begin{frame}[fragile]{Búsqueda en BST: Pseudocódigo}
\begin{verbatim}
Función Buscar(nodo, valor):
    Si nodo es Nulo O nodo.valor == valor:
        Retornar nodo
    
    Si valor < nodo.valor:
        Retornar Buscar(nodo.izquierdo, valor)
    
    Sino:
        Retornar Buscar(nodo.derecho, valor)
\end{verbatim}
\end{frame}

\begin{frame}[fragile]{Estructura del Nodo en Java}
\begin{lstlisting}
class Nodo {
    int valor;
    Nodo izquierdo;
    Nodo derecho;

    public Nodo(int valor) {
        this.valor = valor;
        this.izquierdo = null;
        this.derecho = null;
    }
}
\end{lstlisting}
\end{frame}

\begin{frame}[fragile]{Búsqueda en BST: Java}
\begin{lstlisting}
public Nodo buscar(Nodo raiz, int valorBuscado) {
    if (raiz == null || raiz.valor == valorBuscado) {
        return raiz;
    }

    if (valorBuscado < raiz.valor) {
        return buscar(raiz.izquierdo, valorBuscado);
    } else {
        return buscar(raiz.derecho, valorBuscado);
    }
}
\end{lstlisting}
\end{frame}

\end{document}
